\chapter{System Models}


\section{Digital Communication System}

A general digital communication system transfers information from a source to a destination over a physical transmission medium, referred to as the channel~\cite{syllabus_communicatietheorie}. The corresponding block diagram is shown in Fig.~\ref{fig:digcom_system_block_diagram}.

\begin{figure}[hb]
    \centering
    \makebox[\textwidth][c]{\resizebox{1.05\textwidth}{!}{\includestandalone{figures/part1_background/digcom_system_block_diagram}}}
    \caption{Block diagram of a general digital communication system.}
    \label{fig:digcom_system_block_diagram}
\end{figure}

At the \ac{TX}, the information sequence is first processed by a source encoder to remove redundancy, followed by a channel encoder that introduces controlled redundancy to enable error correction.
The resulting coded bit sequence is partitioned into blocks of $m_c$ bits, each of which is mapped onto a complex-valued symbol drawn from a constellation $\mathcal{X}$ of size $2^{m_c}$.
The modulator then converts the symbol sequence into a baseband signal suitable for transmission.

During transmission, the channel distorts the signal due to attenuation, multipath propagation, and additive noise, such that the received signal differs from the transmitted signal.

At the \ac{RX}, the demodulator converts the received signal into a sequence of decision variables.
These are processed by the detector to obtain estimates of the originally transmitted symbols, which are then mapped back to coded bits.
Finally, the channel and source decoders reconstruct the original information sequence.

In this work, we consider a simplified version of this general architecture in order to focus on the aspects most relevant to our study.
Source coding and decoding are omitted, and the information is assumed to be available as a sequence of \ac{i.i.d.} bits without prior compression or digitization.
Channel coding and decoding are also not considered, such that the analysis is restricted to uncoded transmission.

The transmission medium is assumed to be a wireless channel affected by fading and additive noise. Furthermore, the channel is considered to be frequency-flat, meaning that its response remains constant over the entire signal bandwidth. The specific channel models used in this work are introduced in the relevant chapters.


\section{MIMO}

\subsection{SU-MIMO}

A point-to-point communication system that consists of a \ac{TX} equipped with $N_T$ antennas communicating with a single \ac{RX} equipped with $N_R$ antennas, is called a \ac{SU} \ac{MIMO} system.
A block diagram of such a system is shown in Fig.%~\ref{fig:SU_MIMO_block_diagram}.

Compared to a single-antenna link, the additional spatial degrees of freedom provided by multiple antennas can be exploited to increase data throughput, improve data reliability, or achieve a trade-off between both.
Spatial multiplexing techniques enable significant capacity gains by transmitting multiple independent data streams simultaneously over different antennas, as exemplified by architectures such as BLAST~\cite{foschini_1998}.
Alternatively, spatial diversity can be achieved through space-time coding schemes, which distribute symbols across both space and time to improve robustness against fading. A well-known example is the Alamouti scheme~\cite{alamouti_1998}.
In the future of this work, we focus on spatial multiplexing and in particular on how the we can transmit the data streams simultaneously over different antennas most efficiently.

